%=========== ΛΥΣΗ ΘΕΜΑΤΟΣ Γ.12 ===========
\begin{thema}{Γ}
\begin{erwthma}

\item Η συνάρτηση $f$ είναι δεδομένα παραγωγίσιμη (άρα και συνεχής) στο κλειστό διάστημα $[1, e]$.
Εφαρμόζουμε το \textbf{Θεώρημα Μέσης Τιμής (\eng{T.M.T.})} στο διάστημα αυτό.
Σύμφωνα με το θεώρημα, υπάρχει τουλάχιστον ένα $\xi \in (1, e)$ τέτοιο ώστε η παράγωγος στο σημείο αυτό να ισούται με τον λόγο μεταβολής στα άκρα:
\[ f'(\xi) = \frac{f(e) - f(1)}{e - 1} \]
Αντικαθιστώντας τα δεδομένα $f(e) = e$ και $f(1) = 0$:
\[ f'(\xi) = \frac{e - 0}{e - 1} = \boldsymbol{\frac{e}{e - 1}} \]
Το ζητούμενο αποδείχθηκε.

\item Συγκρίνουμε το κλάσμα με τη μονάδα. Αφού $e > 1$, συνεπάγεται ότι ο παρονομαστής $e-1$ είναι θετικός. 
Ισχύει το προφανές:
\[ e > e - 1 \Rightarrow \frac{e}{e - 1} > \frac{e - 1}{e - 1} \Rightarrow \boldsymbol{\frac{e}{e - 1} > 1} \]
Άρα $f'(\xi) > 1$.

Μας δίνεται η υπόθεση ότι το $\min f'(x) = f'(\xi)$. Αυτό σημαίνει ότι για κάθε $x \in (1, e)$ ισχύει $f'(x) \ge f'(\xi)$.
Από την ανισότητα που μόλις αποδείξαμε, προκύπτει ξεκάθαρα ότι:
\[ f'(x) \ge \frac{e}{e-1} > 1 \Rightarrow f'(x) > 1 \text{  για κάθε } x \in (1, e) \]
Γεωμετρικά, επειδή η $f'(x)$ αποτελεί την κλίση της εφαπτομένης ευθείας $\lambda$, \textbf{συμπεραίνουμε ότι οποιαδήποτε εφαπτομένη της γραφικής παράστασης $C_f$ σε οποιοδήποτε σημείο του διαστήματος $(1, e)$ θα έχει κλίση αυστηρά μεγαλύτερη της μονάδας ($\lambda > 1$)} (δηλαδή σχηματίζει γωνία $> 45^\circ$ με τον άξονα $x'x$).

\item Δίνεται πλέον ότι $f(x) \geq \ln x$ για κάθε $x \in [1, e]$.
Ορίζουμε τη βοηθητική συνάρτηση $g(x) = f(x) - \ln x$. 
Από το δεδομένο έπεται ότι $g(x) \geq 0$ για $x \in [1, e]$.
Ελέγχουμε την τιμή στο αριστερό άκρο του πεδίου:
\[ g(1) = f(1) - \ln 1 = 0 - 0 = 0 \]
Κατά συνέπεια, ισχύει $g(x) \geq g(1)$ για κάθε $x \in [1, e]$. Το σημείο $\boldsymbol{x = 1}$ είναι θέση \textbf{ολικού ελαχίστου} της $g$ στο διάστημα αυτό.
Επειδή το $x=1$ είναι \textbf{άκρο} του διαστήματος και όχι εσωτερικό σημείο (άρα δεν εφαρμόζεται το Θεώρημα \eng{Fermat}), ελέγχουμε τη σχετική παράγωγο μέσω του ορισμού (ή της γενικής κλίσης). Αφού είναι ελάχιστο στο αριστερό άκρο, η συνάρτηση προσδοκά να αυξηθεί (ή να μείνει μηδέν). Οπότε η δεξιά παράγωγός της (η οποία ταυτίζεται με την $g'(1)$ λόγω παραγωγισιμότητας της $f$) θα είναι μη αρνητική:
\[ g'(1) \ge 0 \]
Παραγωγίζουμε την $g$: $g'(x) = f'(x) - \frac{1}{x}$.
Επομένως:
\[ g'(1) \geq 0 \Rightarrow f'(1) - \frac{1}{1} \geq 0 \Rightarrow \boldsymbol{f'(1) \geq 1} \]

\item Ζητείται να εξεταστεί εάν υπάρχει $x_0 \in (1,e)$ με $f(x_0) = x_0 - 1$.
Ορίζουμε $K(x) = f(x) - x + 1$. 
\begin{erat}
\textbf{Προσοχή - Μαθηματικό Σφάλμα Εκφώνησης:} Το ερώτημα αυτό (όπως είναι διατυπωμένο) ενέχει λάθος και η ύπαρξη τέτοιου σημείου είναι \textbf{επισφαλής/αδύνατη} υπό τις τρέχουσες συνθήκες της άσκησης. 

\textbf{Απόδειξη:} Αν δεχθούμε το σενάριο του Ερωτήματος 2, ισχύει $\min f'(x) > 1$, δηλαδή $f'(x) > 1$ $\forall x \in (1, e)$.
Τότε η $K(x) = f(x) - x + 1$ έχει παράγωγο $K'(x) = f'(x) - 1 > 0$. Η $K(x)$ είναι γνησίως αύξουσα. 
Στο $x=1$ ισχύει $K(1) = 0$. Αφού είναι γνησίως αύξουσα, για $x > 1$ προκύπτει υποχρεωτικά $K(x) > K(1) \Rightarrow K(x) > 0$.
Επομένως, σε όλο το $(1, e)$ ισχύει $f(x) - x + 1 > 0 \Rightarrow f(x) > x - 1$. Η συνάρτηση βρίσκεται διαρκώς «πάνω» από την ευθεία $y = x-1$, άρα \textbf{είναι αδύνατον να την τέμνει}. Το αντιπαράδειγμα $f(x) = \frac{e}{e-1}(x-1)$ επαληθεύει όλα τα δεδομένα της άσκησης (παραγωγίσιμη, $f(1)=0, f(e)=e$) αλλά δεν ικανοποιεί πουθενά το $f(x)=x-1$.
\end{erat}
\end{erwthma}
\end{thema}
